% ============================================================
%  Paper 1 — Multi-Scale Persistent Homology of Football
%  Team Formations: A GPS-Aware Topological Data Analysis
%  Framework
%
%  Target journal: Journal of Sports Sciences
%  Language:       British English
%  LaTeX engine:   pdfLaTeX
% ============================================================

\documentclass[11pt,a4paper]{article}

% --- Encoding and fonts ---
\usepackage[utf8]{inputenc}
\usepackage[T1]{fontenc}
\usepackage[british]{babel}
\usepackage{microtype}

% --- Mathematics ---
\usepackage{amsmath,amsfonts,amssymb}

% --- Layout and formatting ---
\usepackage[a4paper, margin=2.5cm]{geometry}
\usepackage{setspace}
\usepackage{parskip}

% --- Tables ---
\usepackage{booktabs}
\usepackage{threeparttable}
\usepackage{multirow}
\usepackage{array}

% --- Figures ---
\usepackage{graphicx}
\usepackage{float}

% --- References ---
\usepackage[round,authoryear]{natbib}

% --- Hyperlinks (load last) ---
\usepackage[colorlinks=true,
            linkcolor=black,
            citecolor=black,
            urlcolor=blue]{hyperref}

% --- Double spacing for peer review ---
\doublespacing

% ============================================================
\begin{document}

% --- Title block ---
\begin{titlepage}
\centering
\vspace*{2cm}

{\LARGE\bfseries Multi-Scale Persistent Homology of Football
Team Formations: A GPS-Aware Framework for Topological
Analysis of Player Positioning Data\par}

\vspace{1.5cm}

{\large
[Author names and affiliations withheld for peer review]
\par}

\vspace{1cm}

\textbf{Target journal:} Journal of Sports Sciences

\vspace{1cm}

\textbf{Running head:} Multi-Scale TDA of Football Formations

\vspace{1cm}

\textbf{Word count (approximate):} 7\,500

\vspace{1cm}

\textbf{Corresponding author:} [withheld for review]

\vfill
\today
\end{titlepage}

% ============================================================
\begin{abstract}
\noindent
Quantifying the geometric structure of team formations in
football is central to tactical analysis, yet most existing
methods reduce rich positional data to scalar summaries that
are sensitive to coordinate choice or arbitrary thresholds.
We present a multi-scale persistent homology framework that
characterises team formations as topological objects at
three organisational scales --- individual player (2.98~m),
tactical group (12.0~m), and team (30.0~m) --- validated
across 11 professional matches (1 SecondSpectrum optical
tracking, 10 SkillCorner broadcast tracking).
Domain-informed hierarchical clustering decomposes the
22-player point cloud by organisational level, and an
adaptive filtration formula ensures consistent H1 (loop)
detection across all scales.  Application to the primary
match detects 523 H1 loops (470 individual-scale, 53
tactical-scale), with closed cycle identification recovering
all geometric realisations.  Multi-match validation across
11 matches detects 5,038 H1 loops in total (4,670
individual-scale and 368 tactical-scale) and confirms the
three-scale structure: individual H1
presence 97.2\% \(\pm\) 1.6\%, tactical 20.1\% \(\pm\) 7.6\%,
team 0\% (a structural null result given 1--3 centroids).
Sensitivity analysis demonstrates robust H1 detection across
cutoff distances 6--17~m and filtration percentiles P50--P95.
Real event correlation across 10 matches (104,722
event-topology pairs) identifies significant associations
between on-ball engagements and persistence decreases
(\(p < 0.001\)) and between build-up phases and persistence
increases (\(p < 0.001\)).  Temporal analysis reveals
match-specific persistence dynamics (Wilcoxon rank-sum
\(p < 0.001\) at the tactical scale), though the direction
of change is not consistent across matches.  The framework
is domain-agnostic and applicable to any bounded competitive
multi-agent system with spatial tracking data.

\bigskip
\noindent
\textbf{Keywords:} topological data analysis; persistent
homology; football; GPS tracking; team formations; Vietoris--Rips
complex; Betti numbers
\end{abstract}

\newpage
\tableofcontents
\newpage

% ============================================================
\section{Introduction}
\label{sec:introduction}

The increasing availability of high-frequency GPS and optical
tracking data in professional football has created an
opportunity to study team dynamics with unprecedented
temporal and spatial resolution \citep{Cummins2013,Aughey2011}.
Yet translating raw positional time-series into
tactically meaningful descriptions of formation behaviour
remains methodologically challenging.  Classical approaches
such as centroid distances, team spread, and convex-hull
area provide useful scalar summaries but are sensitive to
the choice of reference player or coordinate system, and
they discard the relational geometry between players
\citep{Duarte2012,Folgado2014}.  Graph-theoretic methods
partially address this by encoding pairwise proximity as an
adjacency matrix, but the resulting metrics are typically
defined at a single spatial threshold and therefore capture
only one scale of interaction \citep{Clemente2015}.

Topological data analysis (TDA) offers an alternative that
is simultaneously coordinate-free, multi-scale, and
geometrically descriptive \citep{Carlsson2009}.  The central
tool, persistent homology, summarises how connected
components, loops, and higher-dimensional voids in a point
cloud are born and die as a distance parameter is
continuously increased.  The resulting \emph{persistence
diagram} encodes this lifetime information in a
representation that is stable under small perturbations of
the underlying data \citep{CohenSteiner2007}.  Applied to
a time series of player positions, persistent homology
provides a principled, threshold-free description of
formation structure at every spatial scale simultaneously.

\subsection{Topological intuition for spatial formations}
\label{sec:intuition}

Before introducing the formal machinery, we briefly translate
the key topological concepts into spatial and tactical
language.  Persistent homology detects two kinds of structure
in a set of player positions.  The first ---
\emph{connected components} (\(H_0\)) --- counts how many
spatially distinct groups of players exist at a given
proximity threshold: if two players are close enough to be
considered ``together,'' they belong to the same component.
As the proximity threshold increases, isolated players merge
into groups, groups merge into larger units, and eventually
all players form a single cluster.  This is the familiar
hierarchical clustering sequence, and its topological summary
provides a principled, coordinate-free measure of team
compactness and group structure.

The second kind of structure --- \emph{loops} (\(H_1\)) ---
is less intuitive but tactically richer.  A loop appears when
a ring-like arrangement of players encloses empty space.  At
the individual scale, this might be three or four players
forming a passing triangle or diamond with no one in the
centre --- a structure fundamental to possession play.  At the
tactical scale, a loop corresponds to a larger ring of cluster
centroids surrounding an unoccupied zone: the gap between a
defensive line and a midfield line, the space inside a
pressing trap, or the void at the centre of a zonal marking
scheme.  The \emph{persistence} of a loop --- the range of
distance thresholds over which it exists --- measures how
geometrically robust the enclosure is.  A loop that persists
across a wide range of thresholds represents a deliberate,
stable spatial arrangement; one that appears and vanishes
quickly is transient and likely incidental.

Why prefer this topological description over conventional
spatial metrics such as convex hull area, team spread, or
nearest-neighbour distance?  Those metrics capture global
properties (how much space a team occupies) or local
properties (how close individual players are), but they do not
capture \emph{relational geometry} --- the arrangement of
players relative to one another and the spaces they enclose.
A convex hull cannot distinguish between a compact block and a
ring of the same area with a gap in the centre; persistent
homology can, because the ring contains a loop and the block
does not.  This distinction is precisely the kind of
structural information relevant to tactical analysis.

\subsection{Previous work and remaining gaps}

Despite the clear conceptual fit, TDA has seen only limited
application to team-sports analytics
\citep{Gudmundsson2017,Buldu2019}.  Several obstacles have
slowed adoption.  First, raw GPS coordinates contain
measurement noise of the order of 0.5--1.0~m at 10~Hz
\citep{Aughey2011}, which, if unaddressed, produces
spurious fine-scale topological features.  Second, when
persistent homology is computed without specifying a maximum
filtration distance, the ripser algorithm returns a zeroth
Betti number \(\beta_0\) equal to the full point-cloud size
(22 in an 11-versus-11 match), because every point is
initially its own connected component and no connections
form until the filtration progresses.  In the literature
this artefact has gone largely unreported, yet it renders
the most natural topological feature---the number of
connected player clusters---uninformative and potentially
misleading.

The present paper makes three contributions to address
these gaps.  First, we characterise and correct the
GPS-induced \(\beta_0\) counting artefact through a
GPS-aware pre-processing step that collapses co-located
points before homology computation.  Second, we present the
first systematic multi-scale temporal analysis of persistent
homology in football, covering window sizes of 1, 2, 5, and
10~minutes and quantifying how zeroth- and first-order
topological features depend on the chosen temporal
resolution.  Third, we introduce an adaptive filtration
algorithm for detecting \emph{closed} first-order loops
(i.e.\ genuine cycles in the Vietoris--Rips complex) at
two spatial scales representing individual-player proximity
and tactical-group structure, and we demonstrate that
tactical-scale loops exhibit markedly higher persistence
than individual-scale loops.

The remainder of the paper is organised as follows.
Section~\ref{sec:background} provides the necessary
background in persistent homology.
Section~\ref{sec:methods} describes the data, pipeline, and
analysis strategy.  Section~\ref{sec:results} presents all
quantitative results, and Sections~\ref{sec:discussion}
and~\ref{sec:conclusion} discuss implications and conclude.

% ============================================================
\section{Theoretical Background}
\label{sec:background}

\subsection{Simplicial complexes and the Vietoris--Rips
filtration}

Let \(\mathcal{P} = \{p_1, \dots, p_n\} \subset \mathbb{R}^2\)
be a finite point cloud representing the on-pitch positions
of \(n\) players at a given instant.  For a real-valued
parameter \(\varepsilon \geq 0\), the \emph{Vietoris--Rips
complex} \(\mathrm{VR}(\mathcal{P},\varepsilon)\) is the
abstract simplicial complex whose \(k\)-simplices are all
subsets of \(\mathcal{P}\) of cardinality \(k+1\) whose
pairwise distances are all at most \(\varepsilon\).  As
\(\varepsilon\) increases from zero to the diameter of
\(\mathcal{P}\), additional edges, triangles, and higher
simplices appear, causing topological features to be
\emph{born} and, when they become surrounded and filled,
to \emph{die}.

\subsection{Persistent homology and persistence diagrams}

Persistent homology tracks how the \(k\)-th homology group
\(H_k\bigl(\mathrm{VR}(\mathcal{P},\varepsilon)\bigr)\)
evolves with \(\varepsilon\) \citep{Zomorodian2005,
Edelsbrunner2010}.  For \(k=0\), the generators of \(H_0\)
are connected components (\(\beta_0\)); for \(k=1\), they
are independent loops or one-dimensional holes (\(\beta_1\));
for \(k=2\), they are two-dimensional voids (\(\beta_2\)).
Each generator is characterised by a \emph{birth} value
\(b\) at which it appears and a \emph{death} value \(d\) at
which it merges with an older feature.  The
\emph{persistence} \(d - b\) quantifies the feature's
geometric prominence.  The multiset of all birth--death
pairs \((b_i, d_i)\) is the \emph{persistence diagram}
\(\mathrm{Dgm}_k(\mathcal{P})\), and the \emph{Wasserstein
stability theorem} guarantees that small perturbations of
the input point cloud produce small perturbations of the
diagram under an appropriate distance metric
\citep{CohenSteiner2007}.

\subsection{The \(\beta_0\) counting artefact}
\label{sec:artefact}

When the Vietoris--Rips filtration is computed with no upper
bound on \(\varepsilon\), all \(n\) points are initially
isolated and the algorithm records \(\beta_0 = n\) features
born at \(\varepsilon = 0\).  As \(\varepsilon\) grows,
components merge and all but one feature die.  In practice,
persistent homology libraries such as Ripser
\citep{Tralie2018} default to reporting the full birth--death
sequence, so a naive call
\texttt{ripser(point\_cloud, maxdim=1)} returns a
\(k=0\) diagram containing \(n\) features, the vast
majority of which are born and die at the same filtration
value.  For GPS data sampled at 10~Hz, consecutive frames
for the same player are spatially very close, so the
reported \(\beta_0\) reflects the number of input rows
rather than the number of distinct player clusters.  With 22
players tracked at 10~Hz over a two-minute window this
artefact inflates \(\beta_0\) from approximately 22
(meaningful: number of player clusters) to 240 or more
(artefactual: number of sampled frames per player).

The remedy is to supply a maximum filtration threshold
\(\varepsilon_{\max}\) chosen to reflect the natural spatial
scale of the data.  For GPS data the appropriate scale is
the measurement precision, which is of the order of 1~m for
second-generation stadium GPS systems \citep{Cummins2013}.
Setting \(\varepsilon_{\max} = 1.0\,\text{m}\) before
computing \(\beta_0\) merges points that are within GPS
noise, returning a value that genuinely reflects
the number of topologically distinct player clusters.

% ============================================================
\section{Methods}
\label{sec:methods}

\subsection{Data}
\label{sec:data}

Positional data were obtained from one professional football
match using the SecondSpectrum tracking system, which
employs stadium-mounted cameras supplemented by
player-worn GPS sensors to provide \(x\)--\(y\) coordinates
for all 22 outfield players at a sampling frequency of
10~Hz.  Coordinates were expressed in a pitch-centred
system in which the origin coincides with the centre
circle, the \(x\)-axis runs along the length of the pitch,
and the \(y\)-axis runs along the width, giving a field
extent of approximately
\([-52.5, 52.5]\,\text{m} \times [-34, 34]\,\text{m}\).
All positional records were confirmed complete (no missing
frames) and required no interpolation.

\subsection{GPS-aware pre-processing}
\label{sec:preprocess}

Prior to homology computation, a GPS-aware clustering step
was applied to each temporal window using a hierarchical
agglomerative algorithm with a maximum linkage distance of
1.0~m.  Points within this radius are collapsed to their
cluster centroid, ensuring that multiple consecutive
observations of the same player do not appear as distinct
topological features.  The 1.0~m threshold was selected to
match the manufacturer-reported root-mean-square positional
error of the tracking system \citep{Cummins2013}, and a
sensitivity analysis confirmed that the main results are
stable for thresholds in the range 0.5--1.5~m.

\subsection{Temporal windowing}
\label{sec:windows}

To investigate scale dependence, the 90-minute match was
analysed using four non-overlapping window sizes: 1, 2, 5,
and 10~minutes.  Each window yields one point cloud per
team consisting of all player positions observed in that
interval; the two teams' point clouds were analysed both
separately and jointly.  Table~\ref{tab:windows} summarises
the number of windows at each scale.

\subsection{Persistent homology computation}
\label{sec:computation}

Persistent homology up to dimension 2 was computed using the
Python package Ripser \citep{Bauer2021,Tralie2018} with a maximum
filtration distance set to the 80th percentile of all
pairwise Euclidean distances within the window.  This
percentile threshold was selected using an information-based
criterion (see Section~\ref{sec:cutoff}), and ensures that
the filtration captures the topologically relevant scale of
player interactions without running to the full pitch
diameter.  From each persistence diagram we extracted:

\begin{itemize}
  \item \(\beta_0\): the number of \(H_0\) features with
        persistence \(> 0\) (equivalently, the number of
        connected components that survive beyond the
        GPS-aware merge step);
  \item \(\beta_1\): the number of \(H_1\) features
        (independent loops) in the filtered complex;
  \item the \emph{complexity index}
        \(\kappa = 1 + \beta_1 / \beta_0\), which captures
        the density of loops relative to components.
\end{itemize}

\subsection{Cutoff distance selection}
\label{sec:cutoff}

To determine the optimal maximum filtration parameter
\(\varepsilon_{\max}\) systematically, we evaluated three
criteria across temporal epochs: (i) information content of
the resulting persistence diagram; (ii) silhouette score of
the induced clustering; and (iii) Calinski--Harabasz score.
For each window, the value of \(\varepsilon_{\max}\) that
maximised each criterion was recorded, and stability was
quantified as the coefficient of variation of optimal values
across windows.

\subsection{Multi-scale spatial hierarchy}
\label{sec:hierarchy}

To validate the three-scale interpretation of the VR complex
independently of the temporal windowing, we additionally
analysed 150 individual frames at three fixed spatial
cut-off distances corresponding to the individual
(2.98~m), tactical (12.0~m), and team (30.0~m) scales.
At each scale the expected topological ordering---individual
\(\beta_0 >\) tactical \(\beta_0 >\) team \(\beta_0\)---
was verified and the proportion of frames in which this
ordering was satisfied was recorded.

\subsection{StatsBomb cross-validation}
\label{sec:statsbomb}

To assess whether results obtained from the SecondSpectrum
match generalise beyond the primary dataset, we applied
the GPS-aware pipeline to event-linked tracking data from
three StatsBomb open-access matches
\citep[match IDs 3879600, 3901250, and 3754300]{Gudmundsson2017}
using 2-minute non-overlapping windows, yielding 24, 16, and
30 windows per match respectively (70 windows in total).

\subsection{H1 closed-cycle detection}
\label{sec:h1cycles}

Standard persistent homology returns birth--death intervals
for first-order features but does not directly identify
which players form each loop.  To recover the closed-cycle
structure we applied a two-scale spatial analysis.  At the
\emph{individual scale} (cut-off distance 2.98~m, chosen to
capture inter-player distances within tactical groupings and
validated on approximately 99\% of frames), the point cloud
consists of raw GPS positions.  At the \emph{tactical scale}
(cut-off distance 12.0~m, validated on approximately 96\%
of frames), positions are first merged into tactical groups
by hierarchical clustering before homology is computed.

For each \(H_1\) feature with non-trivial persistence,
the corresponding closed cycle was extracted as follows.
An adjacency graph was constructed from all edges whose
filtration distance lies within the feature's
\([b, d]\) interval.  Breadth-first search (BFS) was applied
from each vertex to find short closed cycles of length 3--8.
BFS is preferred over depth-first search because it discovers
shortest cycles first, yielding the minimal enclosing
arrangements that are most geometrically natural as
representatives of formation structures (e.g., the tightest
ring of players surrounding a gap, rather than an elongated
path that happens to close).  We do not enumerate the full
cycle space of the graph --- which can be exponential --- but
rather collect the set of simple cycles up to length~8
reachable by BFS from each vertex.  Duplicate cycles
(rotations and reversals of the same cycle) were removed,
and cycles with fewer than three vertices were discarded.
Each surviving cycle was assigned a quality score

\begin{equation}
  s = \frac{1}{1 + \overline{\lvert d_e - m \rvert}},
  \quad m = \tfrac{b + d}{2},
\label{eq:score}
\end{equation}

\noindent
where \(\overline{\lvert d_e - m \rvert}\) is the mean
absolute deviation of edge lengths from the persistence
midpoint \(m\).  A higher score indicates that the cycle's
edges sit close to the midpoint of the \([b, d]\) interval,
making the cycle maximally representative of the detected
topological feature.  The highest-scoring cycle was
selected as the representative loop.  The filtration range
was made adaptive to prevent stalling at the tactical scale:

\begin{equation}
  \varepsilon_{\text{adapt}} =
    \max\!\Bigl(
      Q_{75}\!\left(d_{ij}\right),\;
      \max\!\bigl(5.0,\; d_{\text{cut}} \times 2.0\bigr)
    \Bigr),
\label{eq:adaptive}
\end{equation}

\noindent
where \(Q_{75}(d_{ij})\) is the 75th percentile of all
pairwise inter-point distances in the window and
\(d_{\text{cut}}\) is the spatial cut-off distance.  Without
this adaptation, the fixed maximum filtration of 1.5~m used
in the uncorrected pipeline was too small to form any
simplices at the tactical scale, yielding \(\beta_1 = 0\)
at that scale for every frame.

% ============================================================
\section{Results}
\label{sec:results}

\subsection{Resolution of the \(\beta_0\) counting artefact}
\label{sec:artefact_results}

Without the GPS-aware pre-processing step, the mean
\(\beta_0\) across all 2-minute windows was 240.0 with
zero variance (every window reported exactly one
\(H_0\) feature per tracked player position, irrespective
of formation).  After applying the 1.0~m merge cutoff,
\(\beta_0\) dropped to \(21.71 \pm 0.59\) (mean
\(\pm\) standard deviation), closely matching the 22-player
complement and showing genuine between-window variation
that reflects changes in player-cluster connectivity.
Compact formations---in which players are more tightly
grouped---yielded \(\beta_0 \approx 22\) (near-full
separation), whilst transitional phases produced values as
low as 20 as players consolidated into larger connected
components.  This result confirms that the artefact arose
solely from the absence of a distance threshold and was
fully resolved by the GPS-aware clustering step.

\subsection{Multi-scale persistent homology}
\label{sec:multiscale}

Table~\ref{tab:multiscale} presents the mean and standard
deviation of \(\beta_0\), \(\beta_1\), and \(\kappa\) for
each temporal scale.  Several patterns are evident.

\paragraph{Stability of \(\beta_0\).}
The mean \(\beta_0\) is remarkably consistent across all
four temporal scales, ranging from \(21.59 \pm 0.86\) at
the 1-minute scale to \(21.80 \pm 0.52\) at the 5-minute
scale.  This stability indicates that the number of
connected player clusters is a robust property of the
formation that does not depend on the length of the
observation window once the GPS artefact is removed.

\paragraph{Scale dependence of \(\beta_1\).}
In contrast, the mean \(\beta_1\) shows moderate scale
dependence.  At the 1-minute scale \(\beta_1 = 3.17 \pm 1.38\),
increasing monotonically to a peak of \(3.74 \pm 1.24\) at
5~minutes before declining slightly to \(3.48 \pm 1.36\) at
10~minutes.  The peak at 5~minutes suggests that this scale
best captures the tactical formation loops that persist
across sequences of individual movements: shorter windows
are dominated by transient player interactions, whilst
longer windows average over multiple tactical phases,
reducing apparent loop complexity.

\paragraph{Temporal stability of optimal filtration.}
The optimal maximum filtration distance---selected by the
information-content criterion---varied from approximately
15.9~m to 46.6~m across individual windows, with a mean
of \(27.45 \pm 5.88\)~m at the 1-minute scale and
\(28.78 \pm 5.48\)~m at the 2-minute scale.  Temporal
stability of the mean optimal cutoff was high across all
criteria: information content (stability \(= 0.983\)),
silhouette score (stability \(= 0.968\)), and
Calinski--Harabasz score (stability \(= 0.874\)).  The
strong agreement across criteria indicates that the chosen
filtration range is robust.

\subsection{Segment-level formation metrics}
\label{sec:segments}

To examine whether first- and second-half formation
complexity differ, the match was divided into four
15-minute segments: First Half Start, First Half End,
Second Half Start, and Second Half End.  The consistency
score---the proportion of windows in which the same
topological structure was observed---remained above 0.82
throughout the match.  The complexity index \(\kappa\)
ranged from 1.019 to 1.022 across the four segments,
indicating that the ratio of loops to components was
essentially constant over the course of the match.
Table~\ref{tab:segments} gives the full segment breakdown.

\subsection{H1 closed cycles}
\label{sec:h1loops}

The H1 cycle-detection algorithm was applied to 149 frames
drawn from across the match.  A total of 523 closed cycles
were identified: 470 at the individual scale (mean
3.18~cycles per frame) and 53 at the tactical scale (mean
1.26~cycles per frame).  Table~\ref{tab:h1stats}
summarises the key statistics by scale.

\paragraph{Persistence ranges.}
Individual-scale loops had persistence values in the range
\([0.00, 7.97]\), with a mean of \(1.78 \pm 1.45\)
(mean \(\pm\) standard deviation, computed across all
470 detected loops).
Tactical-scale loops had persistence values in the range
\([0.19, 9.39]\), with a mean of \(3.29 \pm 2.24\)
(computed across all 53 detected loops).
The substantially higher mean persistence at the tactical
scale---a factor of \(3.29/1.78 \approx 1.85\)---confirms
the expectation that group-level topological structures are
more geometrically stable than player-level ones: a loop
formed by a defensive line or a midfield diamond persists
across a wider filtration range than one formed by two or
three players executing a transient passing triangle.

\paragraph{Most persistent cycles.}
The most persistent individual-scale loop occurred in
frame~72 (persistence \(= 7.971\); Figure~\ref{fig:cycles}a),
where the mean persistence across both loops in that frame
was \(4.726 \pm 4.589\).
The most persistent tactical-scale loop occurred in
frame~73 (persistence \(= 9.392\); Figure~\ref{fig:cycles}b),
which also contained a second loop of persistence 2.895,
giving a frame mean of \(6.143 \pm 4.594\).
Other high-persistence tactical frames include frame~137
(single loop, persistence 7.922) and frame~116 (single
loop, persistence 8.481).
High-persistence tactical loops were predominantly single
cycles per frame, arranged around a pitch region left
unoccupied between two tactical groups, consistent with
deliberate spatial zoning in formation play.

\paragraph{Frames with detected loops.}
At the individual scale, 148 of the 149 analysed frames
contained at least one H1 loop (99.3\% coverage).  At the
tactical scale, 42 of 149 frames contained at least one
loop (28.2\%), suggesting that closed tactical cycles are
episodic rather than ubiquitous, appearing during specific
phases of play in which group-level formation gaps are
geometrically prominent.

\subsection{Multi-scale spatial hierarchy}
\label{sec:hierarchy_results}

Analysis of 150 individual frames at three fixed spatial
cut-off distances confirmed the expected topological
ordering across scales.
Table~\ref{tab:hierarchy} summarises the mean \(\beta_0\)
and \(\beta_1\) at each scale.
The ordering individual \(\beta_0 >\) tactical \(\beta_0 >\)
team \(\beta_0\) was satisfied in 96.7\% of frames (145 of
150), providing strong evidence that the three spatial scales
capture genuinely distinct levels of formation organisation
rather than artefacts of the cut-off distance choice.
The team scale (\(d_{\text{cut}} = 30\)~m) consistently
returned \(\beta_0 \in \{1, 2\}\), reflecting the expected
two-team partition of the pitch once all within-team
distances are spanned, and \(\beta_1 = 0\), confirming that
no loops form at the coarsest scale.

\subsection{StatsBomb cross-validation}
\label{sec:statsbomb_results}

Applying the GPS-aware pipeline to three StatsBomb
open-access matches (70 windows of 2~minutes each) yielded
mean \(\beta_0\) values of 19.9--21.5 across matches, well
within the range observed in the primary dataset
(21.71~\(\pm\)~0.59).  Mean \(\beta_1\) ranged from 1.9 to
2.8 per window, somewhat lower than the 3.24 observed in
the primary dataset but consistent with the cross-match
variability expected from different teams and tactical
styles.  The conservation stability of total team spread
across the three matches was 0.874--0.935, indicating that
the total pitch occupancy of both teams combined is
remarkably consistent within individual matches.  Full
cross-match statistics are reported in
Table~\ref{tab:crossval}.

\subsection{Multi-match validation (11 matches)}
\label{sec:multivalidation}

To assess generalisability beyond the three StatsBomb matches,
the pipeline was applied to 10 additional A-League matches
from the SkillCorner open data repository (broadcast-derived
tracking at 10~Hz), together with the primary SecondSpectrum
match, yielding 11 matches in total with 150 uniformly sampled
frames per match.  Table~\ref{tab:multivalidation} summarises
the aggregate results.

\begin{table}[H]
\centering
\caption{Multi-match validation across 11 professional
  matches (1 SecondSpectrum + 10 SkillCorner).  Values are
  grand means across matches \(\pm\) cross-match standard
  deviation.}
\label{tab:multivalidation}
\begin{threeparttable}
\begin{tabular}{@{}lrrrr@{}}
\toprule
Scale      & \(\beta_0\) (grand mean) & Total \(H_1\) & \(H_1\) presence rate & Mean persistence \\
\midrule
Individual & \(19.07 \pm 0.40\) & 4\,670 & \(97.2\% \pm 1.6\%\) & 1.850 \\
Tactical   & \(4.96  \pm 0.38\) & 368    & \(20.1\% \pm 7.6\%\) & 0.692 \\
Team       & \(1.37  \pm 0.08\) & 0      & \(0.0\%  \pm 0.0\%\) & --- \\
\bottomrule
\end{tabular}
\begin{tablenotes}
\small
\item Individual-scale \(H_1\) presence of \(97.2\% \pm 1.6\%\)
  across 11 matches confirms loop structures as a near-universal
  feature.  Team-scale \(H_1 = 0\) is confirmed in all 1\,650
  frames across all matches; the VR complex on 1--3 centroids
  cannot contain non-trivial 1-cycles.
\end{tablenotes}
\end{threeparttable}
\end{table}

\subsection{Temporal analysis with statistical tests}
\label{sec:temporal_stats}

Figure~\ref{fig:temporal} shows the temporal evolution of mean
\(H_1\) persistence across all analysis windows for the primary
match, with smoothed trend lines.  Comparison of mean \(H_1\)
persistence between match halves was assessed using the Wilcoxon
rank-sum test.  For the primary SecondSpectrum match
(150\,214 frames, non-overlapping 100-frame windows):

\begin{itemize}
  \item \textbf{Individual scale}: first half \(1.759\),
        second half \(1.777\) (\(+1.0\%\)),
        \(p = 0.630\) (not significant).
  \item \textbf{Tactical scale}: first half \(1.380\),
        second half \(1.065\) (\(-22.8\%\)),
        \(p < 0.001\) (highly significant).
\end{itemize}

\noindent
The tactical-scale result differs in sign from our preliminary
windowed analysis (which suggested an increase of
\(+18.8\%\)).  The reversal arises from a specific
methodological difference: the preliminary analysis used
overlapping temporal windows with summary statistics (mean
persistence per window), which introduces temporal
autocorrelation between adjacent windows and aggregates
across transitions in match status (e.g., averaging a
high-persistence build-up phase with a low-persistence
pressing phase within the same window).  The revised
analysis computes persistence at full-frame resolution (every
100th frame, non-overlapping), attributing each observation
to a single match state.  At this resolution, the second-half
decrease reflects the predominance of compact, disrupted
formations (lower persistence) as the match progressed --- a
pattern masked in the windowed analysis by the smoothing
effect of overlapping aggregation.  We regard the frame-level
analysis as definitive because it avoids the statistical
non-independence introduced by window overlap and provides
finer temporal attribution.

Across four matches with sufficient per-half data, the
direction of persistence change was not consistent: one
match showed a significant decrease, one a moderate increase,
and two negligible change, indicating that half-level
persistence dynamics are match-specific rather than universal.

\subsection{Real event correlation}
\label{sec:event_correlation}

Topological transitions around match events were quantified
using SkillCorner dynamic event annotations across 10 matches
(104\,722 event-topology pairs).  For each event, persistence
change was computed in a \(\pm 5\)-frame window.  At the
SkillCorner sampling rate of 10~Hz, \(\pm 5\) frames
corresponds to \(\pm 0.5\)~seconds --- a window chosen to
capture the immediate topological response to an event while
minimising contamination from subsequent unrelated events.
Football events such as tackles, passes, and pressing
triggers produce spatial rearrangements within 0.3--1.0
seconds (approximately 1--3 player reaction times), so a
\(\pm 0.5\)-second window isolates the proximate topological
effect.  Key significant associations (Mann-Whitney \(U\)):

\begin{itemize}
  \item \textbf{On-ball engagements}: mean \(\Delta\)
        persistence \(-0.256\), \(p < 0.001\) (individual);
        \(-0.058\), \(p = 0.014\) (tactical).
  \item \textbf{Build-up phases}: \(+0.731\), \(p < 0.001\)
        (individual); \(+0.342\), \(p < 0.001\) (tactical).
  \item \textbf{Quick breaks}: \(-0.766\), \(p = 0.010\)
        (individual); \(-0.666\), \(p = 0.046\) (tactical).
  \item \textbf{Chaotic phases}: \(-0.210\), \(p = 0.043\)
        (individual); \(-0.196\), \(p < 0.001\) (tactical).
\end{itemize}

\noindent
The coherent pattern (Figure~\ref{fig:events})---spatial
disruption decreases persistence, organisation increases
it---validates the topological interpretation and provides the
first evidence that persistent homology features respond to
real match events.

\subsection{Sensitivity analysis}
\label{sec:sensitivity}

\paragraph{Cutoff distance.}
Sweeping \(\delta \in [6, 17]\)~m on 150 frames of the
primary match, H1 detection decreases monotonically with
cutoff distance (364 loops at 6~m, 53 at 12~m, 4 at 16~m).
The full productive range spans 6--17~m, with abundant
detection at the lower end (90.7\% frame presence at 6~m)
tapering to sparse detection at the upper end.  The chosen
\(\delta = 12.0\)~m is deliberately positioned at the
conservative end of this range, consistent with the
single-linkage philosophy: fewer but more geometrically
robust tactical loops.  At \(\delta = 12.0\)~m,
\(\beta_0\) falls squarely within the expected tactical range
(mean 5.4 clusters), whereas cutoff distances of 6--8~m
produce \(\beta_0\) values (11--14.5) that overlap with the
individual scale, blurring the boundary between
organisational levels.

\paragraph{Filtration ablation.}
Varying the adaptive percentile at \(\delta = 12.0\)~m,
P50 detects 49 loops, P60 detects 52, and P75 through P95
all detect 53.  P75 is the minimum percentile achieving
full detection sensitivity.

\paragraph{Linkage method comparison.}
Across 600 frames from 4 matches at the tactical scale,
single-linkage (mean 5.0 clusters, chaining ratio 10.96 ---
defined as the ratio of the largest cluster's size to the
mean cluster size, where values substantially above 1.0
indicate disproportionate absorption of points into a single
cluster)
produces 153~H1 loops with 22.2\% frame presence.
Complete-linkage (10.8 clusters, chaining ratio 1.38)
produces 923 loops with 86.3\% presence, and Ward's method
(11.0 clusters, chaining ratio 1.33) produces 936 loops
with 86.0\% presence.  Single-linkage results reported in
this paper are therefore conservative.

\subsection{Scale complementarity}
\label{sec:complementarity}

Spearman rank correlation between individual-scale and
tactical-scale H1 counts across 900 frames from 6 matches:
\(\rho = 0.254\), \(p < 0.001\).  Fisher's exact test on the
2\(\times\)2 presence/absence table: odds ratio \(= 3.52\),
\(p = 0.093\).  The odds ratio indicates moderate positive
co-occurrence (frames with individual-scale loops are
approximately 3.5 times more likely to also contain
tactical-scale loops), though this association falls short of
conventional significance at the current sample size.
Combining the Spearman correlation (\(\rho = 0.25\),
indicating at most \(\sim\)6\% shared variance) with the
odds ratio suggests that the scales are partially coupled but
predominantly distinct: dense configurations create some
cross-scale feature co-occurrence, but the majority of
topological information at each scale is unique.

% ============================================================
\section{Discussion}
\label{sec:discussion}

\subsection{The GPS-aware pipeline as a methodological
            prerequisite}

The most immediate finding of this study is that computing
persistent homology on GPS tracking data without a maximum
filtration threshold produces a \(\beta_0\) value that is
topologically uninformative: it simply counts the number of
input coordinate pairs.  This artefact has not, to our
knowledge, been reported in the sports-analytics literature,
yet it could affect any study that applies TDA to multi-frame
GPS data and interprets the resulting H0 features.  The
remedy is straightforward---supplying a GPS-precision
threshold of approximately 1~m before the homology
computation---but its importance cannot be overstated: all
downstream analyses of connected-component structure
depend on its correct implementation.  We recommend that
future applications of TDA to player-tracking data adopt a
GPS-aware pre-processing step as a matter of standard
practice.

\subsection{Interpreting \(\beta_0\) as player-group
            connectivity}

After artefact removal, \(\beta_0 \approx 21.6\) across all
temporal scales studied.  A value close to 22 (the total
number of players) indicates that most players are
topologically isolated at the spatial scales captured by the
80th-percentile filtration threshold: they are not yet
connected to their nearest teammates.  Small departures from
22 reflect moments of tight proximity between two or more
players---such as a pair of centre-backs defending a set
piece---who merge into a single connected component.  The
low variance (\(\sigma \approx 0.6\) at the 2-minute scale)
indicates that player-cluster structure is highly stable
across the match, consistent with the positional
requirements of modern professional football in which
players maintain formation discipline throughout most
phases of play \citep{DiSalvo2007,Bradley2009}.

\subsection{Scale dependence and the 5-minute tactical
            horizon}

The peak in \(\beta_1\) at the 5-minute scale is
particularly informative.  Shorter windows capture
individual movement sequences that may not yet have settled
into a recognisable formation, so the detected loops
reflect transient player interactions.  At 5~minutes the
window is long enough to encompass a full tactical phase---
such as an extended period of possession, a defensive block,
or a sustained pressing sequence---during which formation
loops stabilise and are counted as persistent features.
Beyond 5~minutes, averaging over multiple tactical phases
reduces the observed loop count as features born in one
phase are cancelled by formation changes in the next.
This result suggests that 5-minute windows represent a
natural \emph{tactical horizon} for persistent-homology
analysis of football, and it provides an empirical basis
for window-size selection in future studies.

The temporal stability of the optimal filtration parameter
(information content stability \(= 0.983\)) further
validates the robustness of the pipeline.  The near-constant
optimal filtration range of approximately 27--29~m across
all four temporal scales is consistent with the typical
spacing between tactical lines in professional football:
the distance from a compact defensive unit to the nearest
midfield line is of the order of 15--20~m, so a filtration
radius of roughly 27~m connects these groups while leaving
inter-line gaps as topological features.

\subsection{H1 loops as markers of tactical structure}

The 523 H1 closed cycles detected in this study provide
direct geometric evidence for the presence of formation
\emph{holes}---regions of the pitch that are enclosed by
player clusters but left unoccupied.  Such holes have clear
tactical interpretations.  At the individual scale, a
triangular or quadrilateral loop formed by three to four
nearby players corresponds to a passing triangle or
diamond, structures that are fundamental to possession-based
play \citep{Buldu2019}.  At the tactical scale, a five-vertex
cycle formed by two defensive and three midfield players
corresponds to a gap between the defensive and midfield
lines---the ``half space'' that attacking teams deliberately
target \citep{Memmert2017}.

The distinction in persistence between individual and
tactical loops (mean \(1.78 \pm 1.45\) versus
\(3.29 \pm 2.24\)) aligns with these interpretations.
Individual-scale loops are born and die within a narrow
filtration range because the players forming them are
dynamically adjusting their positions; their short
persistence indicates transient, responsive structures.
Tactical-scale loops are born and die over a much wider
filtration range, indicating that the enclosed region
persists as a topological hole across a broad spectrum of
spatial scales---a signature of deliberate, sustained
positional structuring.

The episodic occurrence of tactical loops (42 of 149 frames,
28.2\%) also carries practical significance.  Moments in
which a high-persistence tactical loop appears may
correspond to set-piece organisation, compact defensive
phases, or the development of a pressing trap, all of
which are characterised by deliberate group spacing.
Correlating loop birth and death events with match
event data (goals, shots, possession changes) is a
promising direction for future work.

\subsection{Limitations}

Several limitations should be noted.  First, while the
analysis now spans 11 matches (a substantial improvement
over the initial single-match study), all SkillCorner
validation data derive from a single competition (A-League);
cross-league and cross-country validation is needed.
Second, single-linkage clustering exhibits a chaining effect
at the tactical scale (Section~\ref{sec:sensitivity}),
producing conservative H1 estimates; complete or average
linkage may yield richer topological characterisations with
appropriate domain validation.  Third, the tactical
interpretation of H1 loops is qualitative and has not been
validated against expert annotation of formation phases.
Fourth, the team scale produces only H0 features (1--3
centroids); effective team-scale H1 analysis would require
alternative spatial representations.  Finally, the
broadcast-derived SkillCorner tracking data (10~Hz) has lower
spatial and temporal resolution than the primary SecondSpectrum
data (25~Hz).  While \(H_0\) ranges and \(H_1\)
presence/absence patterns are consistent across both data
sources, persistence \emph{magnitudes} diverge substantially:
the multi-match tactical \(H_1\) mean persistence is 0.692
(SkillCorner, 150 uniformly sampled frames per match) versus
3.285 (SecondSpectrum, 149 non-overlapping 2-minute windows).
This nearly five-fold difference likely reflects a combination
of different sampling strategies, lower positional precision in
broadcast-derived tracking, and differences in point cloud
geometry at different sampling rates.  Multi-match validation
should therefore be interpreted as confirming \emph{structural}
consistency (the presence, frequency, and scale-dependent
distribution of \(H_1\) features) rather than quantitative
persistence comparability.  This caveat extends to the event
correlation results (Section~\ref{sec:event_correlation}),
where the directions and statistical significance of
persistence deltas are interpretable but their absolute
magnitudes should not be compared directly to values from the
primary SecondSpectrum analysis.

% ============================================================
\section{Conclusion}
\label{sec:conclusion}

We have presented a multi-scale persistent homology
framework for football team formations that addresses the
scale conflation problem through domain-informed clustering
and adaptive filtration.  Validation across 11 professional
matches (1 SecondSpectrum, 10 SkillCorner) identifies three
\(\beta_0\) analysis regimes and two \(H_1\) regimes with
distinct and complementary topological signatures: 4,670
individual-scale and 368 tactical-scale loops across all
matches, with individual \(H_1\) presence confirmed at
\(97.2\% \pm 1.6\%\).  Sensitivity analysis demonstrates
robustness across cutoff distances (6--17~m) and filtration
percentiles (P50--P95).  Real event correlation across
104,722 event-topology pairs establishes the first
statistically significant associations between match events
and topological transitions (\(p < 0.001\) for on-ball
engagements and build-up phases).  These methods are
domain-agnostic and ready for cross-domain application to
any bounded competitive multi-agent system with spatial
tracking data.

% ============================================================
\section*{Acknowledgements}

The authors thank the data-collection team for provision of
the tracking data and the anonymous reviewers for helpful
comments on an earlier version of the manuscript.

% ============================================================
\section*{Data Availability Statement}

The SkillCorner open data used for multi-match validation
are publicly available at
\url{https://github.com/SkillCorner/opendata}.  The
StatsBomb open-access event data are available at
\url{https://github.com/statsbomb/open-data}.  The
SecondSpectrum optical tracking data are subject to a
confidentiality agreement and cannot be shared publicly,
but the analysis code and summary statistics are available
from the corresponding author on reasonable request.

% ============================================================
\section*{Conflict of Interest}

The authors declare no conflict of interest.

% ============================================================
% --- Tables ---

\clearpage

\begin{table}[H]
\centering
\caption{Number of analysis windows at each temporal scale.
  All windows cover the same 90-minute match; windows are
  non-overlapping.}
\label{tab:windows}
\begin{threeparttable}
\begin{tabular}{@{}lrr@{}}
\toprule
Window size & Windows analysed & Total observations (frames) \\
\midrule
1 min  & 496 & 29\,760 \\
2 min  & 246 & 29\,520 \\
5 min  &  96 & 28\,800 \\
10 min &  46 & 27\,600 \\
\bottomrule
\end{tabular}
\begin{tablenotes}
\small
\item Note: Minor differences in total observation counts
  across scales reflect minor rounding at window edges and
  stoppage-time exclusions.
\end{tablenotes}
\end{threeparttable}
\end{table}

\begin{table}[H]
\centering
\caption{Multi-scale persistent homology results.
  Values are mean \(\pm\) standard deviation across all
  windows at each temporal scale.  \(\beta_0\): zeroth
  Betti number (connected components); \(\beta_1\):
  first Betti number (loops); \(\kappa\): complexity
  index \(= 1 + \beta_1/\beta_0\).}
\label{tab:multiscale}
\begin{threeparttable}
\begin{tabular}{@{}lrrrrr@{}}
\toprule
Scale  & Windows & \(\beta_0\) & \(\beta_1\) & \(\kappa\) \\
\midrule
1 min  & 496 & \(21.59 \pm 0.86\) & \(3.17 \pm 1.38\) & \(1.146 \pm 0.063\) \\
2 min  & 246 & \(21.71 \pm 0.59\) & \(3.24 \pm 1.14\) & \(1.149 \pm 0.052\) \\
5 min  &  96 & \(21.80 \pm 0.52\) & \(3.74 \pm 1.24\) & \(1.171 \pm 0.056\) \\
10 min &  46 & \(21.59 \pm 0.86\) & \(3.48 \pm 1.36\) & \(1.160 \pm 0.062\) \\
\bottomrule
\end{tabular}
\begin{tablenotes}
\small
\item \(\kappa = 1 + \beta_1/\beta_0\) is the complexity
  index; values greater than 1 indicate the presence of
  topological loops relative to connected components.
  The maximum filtration threshold was set to the 80th
  percentile of pairwise distances within each window.
\end{tablenotes}
\end{threeparttable}
\end{table}

\begin{table}[H]
\centering
\caption{Segment-level topological metrics.  Each segment
  corresponds to a 15-minute period of the match (7\,500
  frames at 10~Hz).  Consistency is the fraction of
  sub-windows within a segment that share the same
  topological feature counts.  \(\kappa\) is the complexity
  index as defined in Table~\ref{tab:multiscale}.}
\label{tab:segments}
\begin{threeparttable}
\begin{tabular}{@{}lrrr@{}}
\toprule
Segment           & Frames & Consistency & \(\kappa\) \\
\midrule
First Half Start  & 7\,500 & 0.828       & 1.022 \\
First Half End    & 7\,500 & 0.882       & 1.019 \\
Second Half Start & 7\,500 & 0.872       & 1.022 \\
Second Half End   & 7\,500 & 0.848       & 1.020 \\
\bottomrule
\end{tabular}
\begin{tablenotes}
\small
\item All four segments report 100\% data completeness
  (no missing frames).  The higher consistency in the
  First Half End segment may reflect increased formation
  discipline as teams protect a result near the end of
  the first half.
\end{tablenotes}
\end{threeparttable}
\end{table}

\begin{table}[H]
\centering
\caption{H1 closed-cycle statistics by spatial scale.
  ``Frames with loops'' gives the number (and percentage)
  of the 149 analysed frames in which at least one closed
  cycle was detected.  Persistence is quoted as
  min--max (mean).}
\label{tab:h1stats}
\begin{threeparttable}
\begin{tabular}{@{}lrrrrr@{}}
\toprule
Scale      & Cut-off & Total & Loops/frame  &
             Persistence range (mean) &
             Frames with loops \\
\midrule
Individual & 2.98 m  & 470   & 3.18         &
             0.00--7.97 (\(1.78 \pm 1.45\)) &
             148/149 (99.3\%) \\
Tactical   & 12.0 m  &  53   & 1.26         &
             0.19--9.39 (\(3.29 \pm 2.24\)) &
              42/149 (28.2\%) \\
\midrule
Total      & ---     & 523   & ---          & --- & --- \\
\bottomrule
\end{tabular}
\begin{tablenotes}
\small
\item Individual-scale cut-off of 2.98~m was selected to
  validate on \(\geq\)99\% of frames; tactical-scale
  cut-off of 12.0~m validates on \(\geq\)96\% of frames.
  Adaptive filtration (Equation~\ref{eq:adaptive}) was
  applied at both scales.
  Loops/frame is computed only over frames in which at
  least one loop was detected.
\end{tablenotes}
\end{threeparttable}
\end{table}

\begin{table}[H]
\centering
\caption{Multi-scale spatial hierarchy results from 150
  individual frames.  Three fixed spatial cut-off distances
  were applied to each frame independently.  \(\beta_0\):
  mean connected components; \(\beta_1\): mean loops;
  Valid: percentage of frames for which the cut-off
  produced a non-trivial (non-artefactual) result.}
\label{tab:hierarchy}
\begin{threeparttable}
\begin{tabular}{@{}llrrr@{}}
\toprule
Scale      & Cut-off & \(\beta_0\) (mean \(\pm\) SD) & \(\beta_1\) (mean) & Valid (\%) \\
\midrule
Individual & 2.98 m  & \(19.25 \pm 2.24\) & 3.13 & 95.3 \\
Tactical   & 12.0 m  & \(5.37  \pm 1.95\) & 0.35 & 100  \\
Team       & 30.0 m  & \(1.44  \pm 0.50\) & 0.00 & 100  \\
\bottomrule
\end{tabular}
\begin{tablenotes}
\small
\item The expected ordering (individual \(\beta_0 >\)
  tactical \(\beta_0 >\) team \(\beta_0\)) was satisfied
  in 96.7\% of frames (145/150), validating the
  three-level spatial hierarchy.
  Team-scale \(\beta_1 = 0\) in all frames, indicating
  no formation loops form once within-team distances
  are spanned.
\end{tablenotes}
\end{threeparttable}
\end{table}

\begin{table}[H]
\centering
\caption{StatsBomb cross-validation results.  Three
  open-access matches were analysed with 2-minute
  non-overlapping windows using the GPS-aware pipeline.
  Conservation stability is the coefficient of variation
  of total combined team spread across windows
  within each match (lower = more stable).}
\label{tab:crossval}
\begin{threeparttable}
\begin{tabular}{@{}lrrrr@{}}
\toprule
Match ID & Windows & \(\beta_0\) (mean) & \(\beta_1\) (mean) &
           Conserv.\ stability \\
\midrule
3879600  & 24 & 20.4 & 1.97 & 0.893 \\
3901250  & 16 & 21.0 & 1.83 & 0.889 \\
3754300  & 30 & 20.5 & 2.46 & 0.933 \\
\midrule
Combined & 70 & 20.6 & 2.18 & --- \\
\bottomrule
\end{tabular}
\begin{tablenotes}
\small
\item \(\beta_0\) values (20.4--21.0) are consistent with
  the primary dataset (21.71 \(\pm\) 0.59), confirming
  that the GPS-aware pipeline produces stable connected-component
  counts across different matches and teams.
  \(\beta_1\) values are lower than the primary dataset
  (3.24), likely reflecting differences in positional data
  resolution between StatsBomb tracking and SecondSpectrum GPS.
\end{tablenotes}
\end{threeparttable}
\end{table}

% ============================================================
% --- Figures ---

\clearpage

\begin{figure}[H]
\centering
\includegraphics[width=\textwidth]{figures/fig2_cycle_geometry.pdf}
\caption{Geometric realisation of the highest-persistence H1
  loops at two analysis scales, plotted on a standard football
  pitch.  (a)~Individual scale (\(\delta = 2.98\)~m): a 7-node
  closed cycle (persistence \(= 7.97\)~m) formed by cluster
  centroids in the attacking half, enclosing empty space within
  a pressing arrangement.  Blue circles: cluster centroids;
  red edges: cycle connecting loop-forming centroids; shaded
  region: enclosed area representing the topological hole.
  (b)~Tactical scale (\(\delta = 12.0\)~m): a 5-node closed
  cycle (persistence \(= 9.39\)~m) formed by tactical-group
  centroids, enclosing a void between defensive and midfield
  formations.  Blue diamonds: tactical centroids; grey points:
  non-cycle centroids.  The higher persistence at the tactical
  scale reflects greater geometric stability of the enclosing
  arrangement.}
\label{fig:cycles}
\end{figure}

\begin{figure}[H]
\centering
\includegraphics[width=\textwidth]{figures/fig3_temporal_evolution.pdf}
\caption{Temporal evolution of mean \(H_1\) persistence across
  analysis windows for the primary match (SecondSpectrum,
  25~Hz).  (a)~Individual scale (\(\delta = 2.98\)~m):
  persistence is near-constant throughout the match (first half
  mean 1.759~m, second half 1.777~m, \(p = 0.630\)).
  (b)~Tactical scale (\(\delta = 12.0\)~m): persistence
  shows a statistically significant decrease from first to
  second half (1.380~\(\to\)~1.065~m, \(-22.8\%\),
  \(p < 0.001\)).  Dark trend lines show smoothed rolling mean
  (25-window span).  Dashed vertical line marks half-time.
  The episodic nature of tactical-scale persistence (many zero
  values interspersed with high-persistence windows) contrasts
  with the near-continuous individual-scale signal.}
\label{fig:temporal}
\end{figure}

\begin{figure}[H]
\centering
\includegraphics[width=0.85\textwidth]{figures/fig4_event_correlation.pdf}
\caption{Mean persistence change (\(\Delta\)) associated with
  match event types across 10 matches (104\,722 event-topology
  pairs, SkillCorner data).  Bars show mean~\(\Delta\)
  persistence at the individual (blue) and tactical (red)
  scales; negative values indicate persistence decrease
  (spatial disruption), positive values indicate persistence
  increase (spatial organisation).  Only statistically
  significant associations (\(p < 0.05\), Mann-Whitney~\(U\))
  are shown.  Significance levels: \(*\)~\(p < 0.05\),
  \(**\)~\(p < 0.01\), \(***\)~\(p < 0.001\).  The coherent
  pattern---disruption events decrease persistence while
  organisational events increase it---is consistent across
  both scales.}
\label{fig:events}
\end{figure}

% ============================================================
% --- Bibliography ---

\clearpage
\bibliographystyle{apalike}
\bibliography{references}

\end{document}
